% SECCIÓN 1: FINALIDAD PÚBLICA Y ANTECEDENTES
\section{Finalidad pública y antecedentes}

% SUB-SECCIÓN
\subsection{Institución y usuarios}

	La Escuela Profesional de Ingeniería Informática y de Sistemas (EPIIS) de la Universidad Nacional de San Antonio Abad del Cusco (UNSAAC) funciona en el pabellón de la EPIIS, en sus cuatro niveles. Sus usuarios directos son los estudiantes y los docentes, y el personal administrativo también reporta problemas, por lo que debe incluirse en la validación. El marco que regula a la universidad es la Ley Universitaria \cite{leyuniversitaria}.

% SUB-SECCIÓN
\subsection{Servicio y ámbito}

	El servicio evaluado es el acceso a Internet que sostiene las clases, las prácticas, las evaluaciones, la coordinación académica y los trámites en línea. Su alcance comprende los ambientes del pabellón: aulas, laboratorios, oficinas y secretarías, biblioteca, centro federado, cubículos, auditorio, salas de estudio, salas de docentes y de exposiciones, y pasadizos. El inventario exacto de ambientes todavía debe verificarse contra el plano.

% SUB-SECCIÓN
\subsection{Necesidad observada}

	Se han reportado situaciones que, de confirmarse, afectarían el uso académico de la red. En algunas aulas no aparece la señal Wi-Fi. Estudiantes, docentes y personal administrativo señalan cortes de conexión. En los laboratorios se mencionó que sería necesario configurar la dirección IP manualmente, según una declaración atribuida al Ing. Ademir Portillo. También se habló de lentitud, aunque esa observación se refiere a datos móviles personales y no demuestra lentitud en la red institucional. Ninguna de estas situaciones cuenta todavía con registros fechados, actas ni mediciones, por lo que el problema se mantiene como candidato y debe documentarse por ambiente, fecha, hora, actividad afectada y síntoma.

% SUB-SECCIÓN
\subsection{Finalidad pública}

	La finalidad es que los estudiantes y docentes de la EPIIS puedan usar Internet de forma continua durante sus actividades académicas. Esta finalidad se alinea con la prioridad nacional de garantizar el acceso inclusivo, seguro y de calidad al entorno digital \cite{pcm2023}.

% SUB-SECCIÓN
\subsection{Beneficios esperados}

	Si el diagnóstico confirma el problema y la solución resulta eficaz, un acceso más continuo podría facilitar las clases y las prácticas, y contribuir a una mejor experiencia de aprendizaje. Este beneficio es una expectativa y no un resultado demostrado.

% SUB-SECCIÓN
\subsection{Antecedentes pertinentes}

	La UNSAAC cuenta con una tesis que propone el diseño de una red IP-MPLS para mejorar su red de datos \cite{tesisipmpls}. Este estudio es un antecedente de la red de la universidad, pero no reemplaza el diagnóstico del pabellón de la EPIIS. Por otro lado, el equipo ha planteado de manera anticipada una mejora de la infraestructura de red del pabellón, que incluye puntos de red, acceso Wi-Fi, configuración DHCP en laboratorios y ajustes de capacidad o topología. Según el cuadro del equipo, se declaran 25 puntos en el primer nivel, 49 en el segundo, 219 en el tercero y 125 en el cuarto, es decir, 418 en total según esos subtotales. El detalle del cuarto nivel está incompleto y debe cotejarse con el documento fuente. Estas cifras describen la propuesta, pero no prueban que exista una falla ni que la solución sea necesaria.

% SUB-SECCIÓN
\subsection{Estado del antecedente}

	El problema es un candidato sustentado en declaraciones. Faltan episodios documentados por ambiente, periodo, tarea afectada y síntoma, así como mediciones de la red institucional. El periodo de diagnóstico aún debe delimitarse.
